\documentclass[11pt, a4paper]{article}

\input{bfw-style.tex}

\title{\vspace*{-1.5cm}\Huge\bfseries\color{bfwPrimaryDark}Aufgabenheft Session~1\\[0.4em]
  \large\color{bfwInkSoft}\textnormal{IT-Berufe \& duales System --- Übungen im IHK-Klausurformat}}
\author{\textsf{IT Lernfeld 1 --- Das Unternehmen und die eigene Rolle im Betrieb beschreiben}}
\date{Selbstlernmaterial}

\begin{document}
\maketitle
\thispagestyle{empty}

\begin{merksatz}[title={\bfseries So nutzen Sie dieses Aufgabenheft}]
Bearbeiten Sie alle Aufgaben zunächst \emph{ohne} in die Lösungen zu schauen. Schreiben
Sie Ihre Antworten in vollständigen Sätzen --- die IHK-Prüfung verlangt formulierte
Begründungen, nicht bloße Stichworte. Beachten Sie die \emph{Operatoren} (nennen,
erläutern, beschreiben, begründen, beurteilen) --- sie geben den geforderten
Antwort-Tiefegrad vor. Erst danach öffnen Sie den Lösungsteil ab Seite~\pageref{loesungen}
und vergleichen.
\end{merksatz}

\bigskip

\textbf{Kontext:} Sie beginnen als neue Auszubildende bei der \textbf{JIKU IT-Solutions GmbH},
einem IT-Dienstleistungsunternehmen mit 35 Mitarbeitenden. Die Geschäftsführerin, Frau Kaya,
begrüßt Sie und erklärt den Ablauf der Einführungswoche.

\bigskip

\noindent
\begin{tabular}{|l|l|l|l|l|l|l|l|}
\hline
\textbf{Aufgabe} & 1 & 2 & 3 & 4 & 5 & 6 & \textbf{Gesamt}\\
\hline
\textbf{Punkte} & /10 & /8 & /10 & /8 & /10 & /6 & /52\\
\hline
\end{tabular}

\tableofcontents
\vspace{1em}
\hrule

\section{Aufgaben}

\subsection*{Aufgabe~1 --- Die vier IT-Berufe seit 2020 (10 Punkte)}

\begin{enumerate}[label=(\alph*)]
  \item \textbf{Nennen} Sie die vier IT-Ausbildungsberufe, die seit der Neuordnung 2020 existieren. \hfill (4~P.)

  \vspace{3cm}

  \item \textbf{Begründen} Sie, warum der Fachinformatiker in vier Fachrichtungen unterschieden wird. \hfill (3~P.)

  \vspace{3cm}

  \item Welche zwei Fachrichtungen sind seit 2020 \emph{neu hinzugekommen}? \textbf{Erläutern} Sie kurz, warum gerade diese neu eingeführt wurden. \hfill (3~P.)

  \vspace{3cm}
\end{enumerate}

\subsection*{Aufgabe~2 --- Berufsbild zuordnen (8 Punkte)}

\textbf{Ordnen} Sie die folgenden Tätigkeitsbeschreibungen dem passenden IT-Beruf zu.
Jede Beschreibung passt zu genau einem Beruf.

\textit{(Zur Auswahl: Fachinformatiker/-in Anwendungsentwicklung (FIAE), IT-Systemelektroniker/-in,
Kaufmann/-frau für IT-Systemmanagement, Kaufmann/-frau für Digitalisierungsmanagement)}

\begin{enumerate}[label=(\alph*)]
  \item Analysiert Geschäftsprozesse aus betriebswirtschaftlicher Perspektive und erschließt wirtschaftlichen Nutzen aus der Digitalisierung.

  \noindent\rule{\linewidth}{0.3pt}

  \item Erstellt Angebote für IT-Dienstleistungen, administriert IT-Systeme und berät Kunden zu IT-Lösungen.

  \noindent\rule{\linewidth}{0.3pt}

  \item Führt elektrotechnische Arbeiten im IT-Bereich durch, prüft die elektrische Sicherheit und betreut IT-Infrastrukturen.

  \noindent\rule{\linewidth}{0.3pt}

  \item Konzipiert und entwickelt Softwarelösungen für interne und externe Kunden unter Beachtung von IT-Sicherheit und Datenschutz.

  \noindent\rule{\linewidth}{0.3pt}
\end{enumerate}

\medskip
\textit{(2 Punkte je richtige Zuordnung)}

\subsection*{Aufgabe~3 --- Das duale System (10 Punkte)}

Frau Kaya erklärt Ihnen: „Sie lernen jetzt an zwei Orten gleichzeitig."

\begin{enumerate}[label=(\alph*)]
  \item \textbf{Erklären} Sie, was mit dem „dualen System" gemeint ist und welche zwei Lernorte gemeint sind. \hfill (3~P.)

  \vspace{3cm}

  \item \textbf{Beschreiben} Sie je eine Aufgabe des Ausbildungsbetriebs und der Berufsschule im dualen System. \hfill (4~P.)

  \vspace{3cm}

  \item Welches Gesetz bildet die Rechtsgrundlage der dualen Berufsausbildung in Deutschland? \hfill (1~P.)

  \vspace{1.5cm}

  \item In wie vielen staatlich anerkannten Ausbildungsberufen ist eine duale Berufsausbildung in Deutschland möglich? \hfill (2~P.)

  \vspace{1.5cm}
\end{enumerate}

\subsection*{Aufgabe~4 --- Wahr oder falsch? (8 Punkte)}

\textbf{Kreuzen} Sie an und korrigieren Sie falsche Aussagen in einem Satz.

\begin{center}
\begin{tabular}{p{7.5cm} c c p{4.5cm}}
\toprule
\textbf{Aussage} & \textbf{W} & \textbf{F} & \textbf{Korrektur (wenn falsch)}\\
\midrule
a) Die duale Berufsausbildung dauert grundsätzlich zwei Jahre. & {[~]} & {[~]} & \underline{\hspace{4cm}}\\[1.5em]
b) Die IHK ist für Abschlussprüfungen und Beratung zuständig. & {[~]} & {[~]} & \underline{\hspace{4cm}}\\[1.5em]
c) Ausbilder im Betrieb kann jeder Mitarbeiter ohne Eignungsnachweis werden. & {[~]} & {[~]} & \underline{\hspace{4cm}}\\[1.5em]
d) Die Lernfelder 1 bis 6 sind für alle vier IT-Berufe weitgehend gleich. & {[~]} & {[~]} & \underline{\hspace{4cm}}\\[1em]
\bottomrule
\end{tabular}
\end{center}
\textit{(2 Punkte je Zeile: 1 P. Kreuz korrekt, 1 P. Korrektur soweit nötig)}

\subsection*{Aufgabe~5 --- Fallaufgabe: JIKU IT-Solutions (10 Punkte)}

Frau Kaya berichtet, dass JIKU IT-Solutions einen neuen Auszubildenden zum Fachinformatiker
für Systemintegration eingestellt hat: Leon (20 Jahre, mit Hochschulreife).

\begin{enumerate}[label=(\alph*)]
  \item Leon fragt, ob er aufgrund seiner Hochschulreife die Ausbildungsdauer verkürzen kann. \textbf{Erklären} Sie, welche Möglichkeit es dafür gibt und worauf sie rechtlich beruht. \hfill (3~P.)

  \vspace{3cm}

  \item Ein Mitarbeiter von JIKU soll als Ausbilder eingesetzt werden. Er ist 22 Jahre alt und hat die Ausbildereignungsprüfung vor der IHK bestanden. Darf er ausbilden? \textbf{Begründen} Sie Ihre Antwort. \hfill (4~P.)

  \vspace{3.5cm}

  \item \textbf{Beschreiben} Sie, wann und warum Teil~1 der gestreckten Abschlussprüfung stattfindet und wie er in die Gesamtnote eingeht. \hfill (3~P.)

  \vspace{3cm}
\end{enumerate}

\subsection*{Aufgabe~6 --- Handlungskompetenz (6 Punkte)}

\begin{enumerate}[label=(\alph*)]
  \item \textbf{Erläutern} Sie den Begriff „Handlungskompetenz" in eigenen Worten. \hfill (2~P.)

  \vspace{3cm}

  \item \textbf{Ordnen} Sie folgende Eigenschaften den Kompetenzdimensionen \textit{Fachkompetenz}, \textit{Sozialkompetenz} oder \textit{Selbstkompetenz} zu: Toleranz --- Fachkenntnisse --- Flexibilität --- Kooperationsfähigkeit --- Sorgfalt --- Arbeitsqualität. \hfill (4~P.)

  \vspace{3cm}
\end{enumerate}

\vspace{1em}
\noindent\textbf{Punkte gesamt: 52}

\newpage
\section*{Lösungshinweise für Lehrende}\label{loesungen}
\addcontentsline{toc}{section}{Lösungshinweise für Lehrende}

\subsection*{Lösung Aufgabe~1}
\begin{loesung}
\textbf{(a)} Die vier IT-Ausbildungsberufe seit 2020 sind:
1.~IT-Systemelektroniker/-in, 2.~Kaufmann/-frau für IT-Systemmanagement,
3.~Kaufmann/-frau für Digitalisierungsmanagement, 4.~Fachinformatiker/-in.

\textbf{(b)} Wegen der vielfältigen Aufgaben und Einsatzbereiche in der IT-Branche wird ein
breites Spektrum an Fachkenntnissen benötigt, das sich nicht in einer einzigen
Berufsausbildung abbilden lässt. Daher unterscheidet der Beruf Fachinformatiker vier
Fachrichtungen, die jeweils spezifische Kompetenzen vermitteln.

\textbf{(c)} Neu seit 2020 sind \emph{Digitale Vernetzung} (FIDV) und \emph{Daten- und
Prozessanalyse} (FIDP), weil in diesen Bereichen zukünftig hoher Bedarf erwartet wird
(Vernetzung von Produktionssystemen, datengesteuerte Prozessoptimierung).
\end{loesung}

\subsection*{Lösung Aufgabe~2}
\begin{loesung}
(a) Kaufmann/-frau für \emph{Digitalisierungsmanagement}\\
(b) Kaufmann/-frau für \emph{IT-Systemmanagement}\\
(c) \emph{IT-Systemelektroniker/-in}\\
(d) Fachinformatiker/-in \emph{Anwendungsentwicklung}
\end{loesung}

\subsection*{Lösung Aufgabe~3}
\begin{loesung}
\textbf{(a)} Das duale System bezeichnet die Berufsausbildung an zwei Lernorten
gleichzeitig: dem \emph{Ausbildungsbetrieb} für die praktische Ausbildung und der
\emph{Berufsschule (BBS)} für den theoretischen Unterricht.

\textbf{(b)} Ausbildungsbetrieb: schließt Berufsausbildungsvertrag ab und vermittelt
praktische berufliche Handlungsfähigkeit. Berufsschule: erteilt theoretischen Unterricht
in Lernfeldern und allgemeinbildenden Fächern auf Basis des Rahmenlehrplans.

\textbf{(c)} Das \emph{Berufsbildungsgesetz (BBiG)}.

\textbf{(d)} In über \emph{300} staatlich anerkannten Ausbildungsberufen.
\end{loesung}

\subsection*{Lösung Aufgabe~4}
\begin{loesung}
(a) \textbf{Falsch} --- Die Regelausbildungsdauer beträgt grundsätzlich \emph{drei} Jahre.\\
(b) \textbf{Richtig}\\
(c) \textbf{Falsch} --- Ausbilder muss persönlich und fachlich geeignet und mindestens
24 Jahre alt sein; die IHK stellt die Eignung fest.\\
(d) \textbf{Richtig}
\end{loesung}

\subsection*{Lösung Aufgabe~5}
\begin{loesung}
\textbf{(a)} Eine Verkürzung der Ausbildungszeit ist auf Antrag möglich, wenn anerkannte
Bildungsgänge angerechnet werden (z.\,B.\ Hochschulreife). Rechtsgrundlage sind die
§§~7, 8, 45 BBiG. Leon kann bei der IHK einen entsprechenden Antrag stellen.

\textbf{(b)} Nein, er darf \emph{nicht} ausbilden, auch wenn er die
Ausbildereignungsprüfung bestanden hat. Das Lehrbuch gibt an, dass der Ausbilder
mindestens \textbf{24 Jahre} alt sein muss. Mit 22 Jahren fehlt ihm die erforderliche
Altersvoraussetzung. Die IHK stellt die persönliche und fachliche Eignung fest.

\textbf{(c)} Teil~1 findet nach etwa 15 Monaten (Mitte des zweiten Ausbildungsjahres)
statt. Er prüft den Stoff der Lernfelder 1 bis 6 und ist für alle IT-Berufe gleich
(Leitthema: „Einrichten eines IT-gestützten Arbeitsplatzes", 90 Minuten schriftlich).
Teil~1 geht mit \textbf{20\,\%} in die Gesamtnote der Abschlussprüfung ein.
\end{loesung}

\subsection*{Lösung Aufgabe~6}
\begin{loesung}
\textbf{(a)} Handlungskompetenz ist die Bereitschaft und Befähigung, sich in beruflichen,
privaten und gesellschaftlichen Situationen sachgerecht, individuell, durchdacht und sozial
verantwortlich zu verhalten. Sie umfasst Fach-, Sozial- und Selbstkompetenz.

\textbf{(b)}
\emph{Fachkompetenz}: Fachkenntnisse, Arbeitsqualität\\
\emph{Sozialkompetenz}: Toleranz, Kooperationsfähigkeit\\
\emph{Selbstkompetenz}: Flexibilität, Sorgfalt
\end{loesung}

\end{document}
